% Standalone wrapper so the report sections can be compiled and checked on their
% own, and dropped into a larger document unchanged.
%
%   latexmk -pdf main.tex   (or: pdflatex main && bibtex main && pdflatex main x2)
\documentclass[11pt,a4paper]{article}

\usepackage[utf8]{inputenc}
\usepackage[T1]{fontenc}
\usepackage{amsmath,amssymb}
\usepackage{bm}
\usepackage[margin=2.5cm]{geometry}
\usepackage{hyperref}
\usepackage{booktabs}
\usepackage{natbib}   % \citet in related_work
\usepackage{graphicx}
% Figures are the same PNGs the README uses, regenerated by scripts/ -- there is no
% second copy to drift out of date.
\graphicspath{{../assets/}{../assets/casque/}{../assets/pavillon/}}

\title{Video to 3D Gaussian Splatting:\\Two Captures, and Which Findings Transfer}
\author{Adrien Kinart}
\date{\today}

\begin{document}
\maketitle

\begin{abstract}
We reconstruct a carved wooden ceiling panel from handheld video using 3D Gaussian
Splatting. The capture is adversarial for the method: single-sided, low-parallax
and close-range on a near-planar relief. We describe the theory behind each
component of the pipeline and the four regularisation mechanisms added to adapt it
to this regime --- global structure-from-motion, monocular depth priors, spatial
bounding with floater suppression, and per-image appearance modelling for
multi-clip merging.

We then reconstruct a second, deliberately opposite capture --- a full orbit of a
free-standing helmet --- to separate findings that are properties of the method from
findings that were properties of the first capture. Two of three predictions we carried
over proved wrong, informatively: the optimal Gaussian budget \emph{inverts} between the
two captures, and pose refinement remains harmful even where poses are measurably worse.
Throughout we judge differences by paired per-view intervals rather than aggregate means,
which repeatedly changes the conclusion.
\end{abstract}

% !TeX root = main.tex
%
% Related work for the Pavillon 3DGS report. Organised by the axes on which our
% capture is hard, so each family is judged by whether it addresses OUR failure
% modes rather than by general popularity.

\section{Related Work}
\label{sec:related}

Our capture is adversarial along four specific axes: it is \emph{single-sided}
(the camera never orbits the object), \emph{low-parallax}, \emph{close-range} on a
\emph{near-planar relief}, and it is assembled from handheld clips with
independent auto-exposure. We therefore organise the literature by those axes and,
where we tested a family, report what it did on this data.

%---------------------------------------------------------------------------
\subsection{Primitive-based radiance fields and frameworks}

3DGS \cite{kerbl2023gaussian} replaced NeRF's ray-marched implicit field
\cite{mildenhall2020nerf} with explicit anisotropic primitives rasterised by EWA
splatting \cite{zwicker2001ewa,zwicker2001surface}, trading a continuous
representation for real-time rendering and a far cheaper optimisation. Two open
implementations now dominate practice: \texttt{gsplat} \cite{ye2024gsplat}, a
CUDA-backed library exposing the rasteriser and density-control strategies, and
Nerfstudio \cite{tancik2023nerfstudio}, whose \emph{Splatfacto} model wraps
\texttt{gsplat} with a configurable training pipeline. We build directly on the
packaged \texttt{gsplat} APIs rather than on Nerfstudio: the pipeline needs
per-stage checkpointing, preemption handling and fingerprint-driven reruns on a
Slurm cluster, which is orchestration Nerfstudio does not target, and avoiding the
Nerfstudio dependency also avoids \texttt{tiny-cuda-nn}, a known pain point on
Blackwell (\texttt{sm\_120}).

Anti-aliasing is orthogonal but relevant to close-range capture: Mip-Splatting
\cite{yu2024mipsplatting} adds 3D smoothing and a 2D Mip filter so primitives
observed at sampling rates far from training do not produce erosion or dilation
artefacts. Our capture varies sharply in scale between contextual and extreme
close-up views, so we enable \texttt{gsplat}'s antialiased rasterisation mode
throughout.

%---------------------------------------------------------------------------
\subsection{Density control: from heuristics to sampling}
\label{sec:rw-density}

Vanilla 3DGS grows the model with a hand-tuned adaptive density control (ADC):
clone or split wherever the view-space positional gradient exceeds a threshold,
periodically reset opacities, prune the transparent. The thresholds are heuristic
and, as we found empirically (\S\ref{sec:theory} and our capacity experiments),
their interaction with scene scale is fragile.

Two lines of work replace or restructure this. \textbf{Structured
representations} anchor primitives to a coarser scaffold: Scaffold-GS
\cite{lu2024scaffoldgs} distributes local Gaussians around anchor points and
predicts their attributes on the fly from view direction and distance, which
suppresses the redundant floaters that unconstrained ADC produces; Octree-GS
\cite{ren2024octreegs} extends this with an LOD hierarchy so the level of detail
matches the observation scale.

\textbf{3DGS-MCMC} \cite{kheradmand2024mcmc} is the more fundamental reframing and
the most relevant to our findings. It reinterprets the Gaussians as samples drawn
from a distribution over scenes, turns the optimiser into Stochastic Gradient
Langevin Dynamics by injecting noise, and rewrites cloning and pruning as
\emph{state transitions that approximately preserve sample probability} — replacing
``dead'' primitives by relocating them onto live ones rather than heuristically
splitting. Crucially, it reports better reconstructions \emph{at an identical
Gaussian budget}, and it makes the budget an explicit parameter rather than an
emergent consequence of thresholds.

This connects directly to our own result. We measured a monotone preference for
\emph{fewer} primitives on this capture: reducing \texttt{cap\_max} from 1.5\,M to
375\,k raised test PSNR by 1.8\,dB, improved SSIM, and eliminated the
catastrophic-view tail, while doubling it to 3.0\,M degraded every metric. The
compaction literature — Compact-3DGS \cite{lee2024compact3dgs}, GaussianSpa
\cite{zhang2024gaussianspa} — reports that aggressive pruning preserves or improves
quality while shrinking models by large factors, which is consistent, though that
work is usually motivated by memory rather than generalisation. The sparse-view
literature makes the generalisation argument explicitly: DropGaussian
\cite{park2025dropgaussian} randomly removes primitives during training as a
\emph{structural regulariser} for few-view inputs, and recent analysis quantifies
\emph{co-adaptation} between primitives as a distinct sparse-view failure mode
\cite{coadaptation2025}. Our capacity sweep is an independent, purely empirical
instance of the same phenomenon, obtained without any architectural change: on a
low-parallax capture, primitive count behaves as a regularisation hyper-parameter,
not a quality dial. Adopting 3DGS-MCMC would be the principled next step, since it
makes exactly this budget explicit.

%---------------------------------------------------------------------------
\subsection{Sparse and few-view reconstruction}

When views are few or share little parallax, the photometric objective is
underdetermined: many primitive configurations reproduce the training images while
placing mass at wrong depths. The dominant remedy is to import geometry from a
monocular depth network \cite{ranftl2022midas,yang2024depthanythingv2}, used
scale-invariantly because such predictions are only defined up to an affine
transform \cite{eigen2014depth}. FSGS \cite{zhu2024fsgs} densifies guided by depth
and proximity; SparseGS \cite{xiong2024sparsegs} combines depth with a
diffusion-based score to remove background collapse; DNGaussian
\cite{li2024dngaussian} contributes global-local depth normalisation so that
absolute depth error does not dominate local shape. CoR-GS \cite{zhang2024corgs}
takes a different route, training two models and using their disagreement to
identify unreliable geometry.

We implement the depth-prior family (DepthAnything-v2 with a Pearson-correlation
loss, \S\ref{sec:sparse}) and measure it honestly: on this capture the prior is
\emph{nearly free} in PSNR terms (an ablation with it disabled recovers only
0.26\,dB) and slightly \emph{increases} floater removal — but it is not what
determines final quality. Resolution and capacity dominate it by a wide margin.
This is a useful negative-ish result: depth priors are routinely presented as the
answer to sparse-view degradation, yet on a capture whose views are numerous but
badly distributed, the binding constraints were elsewhere.

A newer direction sidesteps SfM entirely, initialising from learned stereo:
DUSt3R \cite{wang2024dust3r} and VGGT \cite{wang2025vggt} regress dense geometry
and poses directly, and InstantSplat \cite{fan2024instantsplat} couples such an
initialisation to 3DGS for very sparse inputs. For our data the SfM stage was not
the bottleneck once global mapping was used (below), so we did not pursue it;
it remains the natural fallback for captures where COLMAP/GLOMAP fail outright.

%---------------------------------------------------------------------------
\subsection{Camera calibration and pose refinement}

Incremental SfM \cite{schoenberger2016sfm} registers images greedily and can
abandon weakly connected regions — exactly the regime of a low-overlap capture.
Global SfM re-estimates all poses jointly; GLOMAP \cite{pan2024glomap} performs
global rotation averaging followed by joint global positioning. Substituting
GLOMAP for incremental COLMAP on identical features and matches raised
registration from 82 to 181 of 193 images on our data, a $2.2\times$ coverage gain
and the single largest early improvement we measured.

Where poses remain inaccurate, they can be refined jointly with the scene, as in
BARF \cite{lin2021barf}, whose central observation is that such optimisation must
be coarse-to-fine to avoid poor local minima; COLMAP-Free 3DGS
\cite{fu2024colmapfree} pushes this to joint pose-and-scene estimation without SfM
at all. We implemented per-camera SE(3) refinement and found it \emph{harmful}
here ($-1$\,dB): GLOMAP's poses were already sub-pixel (0.92\,px mean reprojection
error), so there was no calibration error to recover, and refining only the
training cameras — held-out poses must not be optimised, or the metric is
contaminated — introduced a small relative misalignment between the reconstructed
scene and the un-refined evaluation cameras. Pose refinement is a remedy for bad
poses, and ours were not bad.

%---------------------------------------------------------------------------
\subsection{Surfaces, meshes and hybrid representations}

A radiance field is not a surface, and our subject is a carved relief whose value
is geometric. Several families address this. \textbf{Flattening the primitive}:
2DGS \cite{huang20242dgs} replaces ellipsoids with oriented disks that are surface
elements by construction, and SuGaR \cite{guedon2024sugar} regularises 3D Gaussians
towards a surface before Poisson meshing. \textbf{Reformulating the depth or
opacity}: Gaussian Opacity Fields \cite{yu2024gof} extracts level sets of an
opacity field with Marching Tetrahedra, avoiding Poisson artefacts in unbounded
scenes; PGSR \cite{chen2024pgsr} compresses Gaussians to planes with multi-view
geometric consistency; RaDe-GS \cite{zhang2024radegs} derives a rasterised depth
that is unbiased for the level set. \textbf{Hybrid supervision}: GSDF
\cite{yu2024gsdf} runs a 3DGS branch and a neural SDF branch that supervise each
other — the splatting branch's rendered depth guides SDF ray sampling, and the SDF
in turn guides density control to grow primitives near the surface and prune those
that drift off it. GSDF is the most attractive of these for our problem, because
its density control is driven by a \emph{geometric} signal rather than a
photometric gradient, which is precisely the signal our capture lacks.

We evaluated 2DGS and report a clear negative result: it collapses to flat,
featureless renders (PSNR $\approx$13) on this capture. The explanation is
structural rather than incidental — a disk must recover an orientation as well as a
position, and normals are constrained mainly by observing a surface from multiple
directions, which a single-sided capture does not supply. The added constraint
removes degrees of freedom the photometric loss was using without supplying the
multi-view evidence needed to fix the rest. We therefore extract geometry by
TSDF-fusing \cite{curless1996tsdf} the depth composited by the \emph{3D} Gaussian
model, accepting the known bias of volumetric primitives, and obtain a usable mesh
of the panel. GSDF or GOF would be the principled improvements.

%---------------------------------------------------------------------------
\subsection{Appearance variation across captures}

Merging clips is the cheapest way to buy coverage, but each handheld clip carries
its own auto-exposure and white-balance response, and 3DGS has no parameter
expressing ``same radiance, different camera response''. NeRF-W
\cite{martinbrualla2021nerfw} introduced per-image appearance latents optimised
jointly with the scene, in the GLO style \cite{bojanowski2018glo}. The idea has
since been ported to splatting: Splatfacto-W \cite{xu2024splatfactow} adds
per-Gaussian neural colour features, per-image appearance embeddings and a
spherical-harmonic background to Nerfstudio's Splatfacto; WildGaussians
\cite{kulhanek2024wildgaussians} combines appearance modelling with DINO features
for occlusion robustness. A lighter-weight alternative now standard in Nerfstudio
is the bilateral grid of \citet{wang2024bilateral}, which learns a smooth
low-resolution 3D colour-transform grid per image — more expressive than a global
affine map while still spatially constrained.

We implement the low-capacity end of this spectrum deliberately: a per-image latent
decoded to a \emph{global} affine colour transform. The capacity argument matters
more than the architecture. A per-pixel correction could memorise each training
view and destroy the multi-view constraint; our transform applies identically at
every pixel, so it commutes with any spatial permutation of the image and therefore
carries exactly zero spatial information. It can absorb exposure, gain and
white-balance drift and provably nothing else. A bilateral grid sits between these
extremes and would be the natural upgrade if per-clip vignetting or spatially
varying response mattered.

Our merge nevertheless underperformed the best single-clip model. Diagnosing it
required first fixing a measurement bug — evaluation must apply the appearance
correction, and must do so using a latent that does not depend on the held-out
image, or the metric is either penalised or contaminated. With per-clip latents the
merged model reaches 22.0\,dB against 24.9\,dB single-clip. The residual gap is not
photometric: the second clip enlarges the reconstructed volume while contributing
few views, and the capacity effect of \S\ref{sec:rw-density} then dominates.

%---------------------------------------------------------------------------
\subsection{Positioning}

Relative to this literature our contribution is not a new representation but a
measured account of \emph{which} interventions matter for one well-characterised
hard capture, with negative results reported alongside positive ones. Global SfM
and resolution dominated; capacity reduction was the largest late gain; depth
priors were nearly free but not decisive; 2DGS, multi-clip merging and pose
refinement each degraded quality for reasons we can attribute. The methods we did
not implement — 3DGS-MCMC for principled budget control, GSDF or GOF for geometry,
bilateral grids for appearance — are exactly the ones this evidence points to.

% !TeX root = main.tex
%
% Theory section for the Pavillon 3D Gaussian Splatting report.
% Compile with:  latexmk -pdf main.tex   (see main.tex in this directory)

\section{Theoretical Background}
\label{sec:theory}

This section develops the theory underlying the reconstruction pipeline: the
Gaussian-splatting scene representation and its differentiable rasteriser
(\S\ref{sec:3dgs}), the structure-from-motion stage that supplies camera poses
(\S\ref{sec:sfm}), the adaptive density control that grows the model
(\S\ref{sec:densification}), and the four regularisation mechanisms that adapt the
method to our capture: monocular depth priors, spatial bounding, floater
suppression, and per-image appearance modelling
(\S\ref{sec:sparse}--\S\ref{sec:appearance}). Throughout, we highlight
\emph{why} each component is required by the specific difficulty of our data: a
single-sided, low-parallax, close-range capture of a planar carved relief.

%---------------------------------------------------------------------------
\subsection{Scene representation: 3D Gaussian Splatting}
\label{sec:3dgs}

Neural radiance fields represent a scene as a continuous function evaluated by
ray marching \cite{mildenhall2020nerf}, which couples rendering cost to sample
count. 3D Gaussian Splatting (3DGS) \cite{kerbl2023gaussian} instead adopts an
explicit, unstructured set of anisotropic Gaussian primitives that are
\emph{rasterised}, giving real-time rendering while remaining fully
differentiable.

The scene is a set $\mathcal{G}=\{g_i\}_{i=1}^{N}$ where each primitive carries a
mean $\bm{\mu}_i\in\mathbb{R}^3$, a covariance $\bm{\Sigma}_i\in\mathbb{R}^{3\times3}$,
an opacity $o_i\in\mathbb{R}$ and view-dependent colour coefficients. The spatial
support is the unnormalised Gaussian
\begin{equation}
  G_i(\bm{x}) \;=\; \exp\!\Big(-\tfrac{1}{2}(\bm{x}-\bm{\mu}_i)^{\!\top}\,
  \bm{\Sigma}_i^{-1}\,(\bm{x}-\bm{\mu}_i)\Big).
  \label{eq:gaussian}
\end{equation}

\paragraph{Parameterisation of the covariance.}
A covariance matrix must remain symmetric positive semi-definite, which gradient
descent on its six free entries does not guarantee. 3DGS therefore factorises it
into a rotation and an axis-aligned scaling,
\begin{equation}
  \bm{\Sigma}_i \;=\; \bm{R}(\bm{q}_i)\,\bm{S}_i\,\bm{S}_i^{\!\top}\,\bm{R}(\bm{q}_i)^{\!\top},
  \qquad
  \bm{S}_i=\operatorname{diag}\!\big(\exp(\bm{s}_i)\big),
  \label{eq:covariance}
\end{equation}
optimising a unit quaternion $\bm{q}_i$ and \emph{log}-scales
$\bm{s}_i\in\mathbb{R}^3$. The exponential keeps scales strictly positive and makes
the parameterisation multiplicative, so a single learning rate behaves sensibly
across the wide range of primitive sizes in a scene. Opacity is likewise stored as
a logit and mapped through a sigmoid, $\alpha_i=\sigma(o_i)\in(0,1)$. These
reparameterisations matter operationally: our floater diagnostics
(\S\ref{sec:floaters}) threshold $\sigma(o_i)$ and $\max\exp(\bm{s}_i)$ rather
than the raw parameters.

\paragraph{Projection and the 2D covariance.}
Rendering requires the image-space footprint of each primitive. Under a viewing
transform $\bm{W}$ (world-to-camera) followed by a projective map, the projection
is nonlinear, so 3DGS adopts the EWA splatting approximation
\cite{zwicker2001ewa,zwicker2001surface}: linearise the projective map about
$\bm{\mu}_i$ with Jacobian $\bm{J}$ and propagate the covariance,
\begin{equation}
  \bm{\Sigma}'_i \;=\; \bm{J}\,\bm{W}\,\bm{\Sigma}_i\,\bm{W}^{\!\top}\bm{J}^{\!\top},
  \label{eq:proj}
\end{equation}
retaining the upper-left $2\times2$ block as the screen-space covariance. The
first-order approximation is accurate when a primitive is small relative to its
depth --- an assumption that our extreme close-up views stress, since a primitive
covering a large solid angle has a genuinely non-Gaussian footprint.

\paragraph{Compositing.}
Colour is accumulated by front-to-back alpha compositing, the discretisation of
the volume-rendering integral under an absorption-emission model
\cite{max1995optical}. For a pixel $\bm{p}$, with primitives sorted by depth,
\begin{equation}
  \bm{C}(\bm{p}) \;=\; \sum_{i=1}^{N} \bm{c}_i(\bm{d})\,\hat{\alpha}_i(\bm{p})
  \prod_{j<i}\big(1-\hat{\alpha}_j(\bm{p})\big),
  \qquad
  \hat{\alpha}_i(\bm{p}) = \alpha_i\,G'_i(\bm{p}),
  \label{eq:compositing}
\end{equation}
where $G'_i$ is the projected 2D Gaussian~\eqref{eq:proj} evaluated at the pixel
and $\bm{c}_i(\bm{d})$ is the colour emitted toward the viewing direction
$\bm{d}$, expanded in spherical harmonics,
$\bm{c}_i(\bm{d})=\sum_{\ell=0}^{L}\sum_{m=-\ell}^{\ell}\bm{k}_{i,\ell m}\,Y_{\ell m}(\bm{d})$.
We use $L=3$ (16 coefficients per colour channel) and, following common practice,
introduce the bands progressively during training so that low-frequency albedo is
fitted before view-dependent effects.

Because \eqref{eq:compositing} is a product of per-primitive terms, every quantity
is differentiable and gradients flow to positions, covariances, opacities and
colours simultaneously. This is also the source of a characteristic failure mode:
the same pixel colour can be explained either by a correctly placed opaque
primitive or by a stack of misplaced semi-transparent ones. Nothing in
\eqref{eq:compositing} prefers the former, which is precisely why explicit
regularisation (\S\ref{sec:floaters}) is necessary.

\paragraph{Photometric objective.}
Training minimises a combination of an $\ell_1$ term and a structural
dissimilarity term \cite{wang2004ssim},
\begin{equation}
  \mathcal{L}_{\text{photo}} \;=\; (1-\lambda)\,\big\lVert \hat{\bm{I}}-\bm{I}\big\rVert_1
  \;+\; \lambda\,\big(1-\mathrm{SSIM}(\hat{\bm{I}},\bm{I})\big),
  \label{eq:photo}
\end{equation}
with $\lambda=0.2$. The SSIM term supplies a spatially structured gradient that a
purely per-pixel loss lacks, which matters for high-frequency relief detail.
Optimisation uses Adam \cite{kingma2015adam} with per-parameter-group learning
rates; the positional learning rate is decayed exponentially so primitives settle
rather than jitter, a detail that proved decisive for sharpness in our runs.

%---------------------------------------------------------------------------
\subsection{Camera calibration: incremental versus global SfM}
\label{sec:sfm}

3DGS optimises geometry \emph{given} camera poses; those come from
structure-from-motion. The incremental pipeline of COLMAP
\cite{schoenberger2016sfm} extracts SIFT features \cite{lowe2004sift}, matches
them, then seeds a two-view reconstruction and repeatedly registers one image at a
time, triangulating and running bundle adjustment. Incremental growth is robust on
well-connected sets but is greedy: registration proceeds only where the current
model already provides sufficient $2$D--$3$D correspondences, so a weakly
connected region can be abandoned entirely.

Global SfM instead estimates all camera poses simultaneously from the epipolar
graph. GLOMAP \cite{pan2024glomap} performs global rotation averaging followed by
a joint global positioning step that estimates camera centres and points together,
avoiding the classical fragility of separate translation averaging. Because it
considers the whole view graph at once, it does not inherit the incremental
pipeline's order dependence.

This distinction was decisive for our capture. On the single-orbit Pavillon clip,
incremental COLMAP registered $82$ images whereas GLOMAP registered $181$ of
$193$ from identical features and matches --- a $2.2\times$ coverage gain. Low
image overlap produces exactly the weakly connected graph that penalises greedy
registration, and coverage is the binding constraint on downstream quality: a
surface observed by no registered camera cannot be reconstructed at all.

%---------------------------------------------------------------------------
\subsection{Adaptive density control}
\label{sec:densification}

A fixed primitive count cannot match a scene's spatially varying complexity, so
3DGS \emph{adapts} the set during optimisation \cite{kerbl2023gaussian}. The
driving signal is the magnitude of the loss gradient with respect to the
\emph{projected} (screen-space) mean, averaged over the views in which a primitive
was visible,
\begin{equation}
  \tau_i \;=\; \Big\lVert \frac{\partial \mathcal{L}}{\partial \bm{\mu}'_i}\Big\rVert.
\end{equation}
A large $\tau_i$ indicates that the renderer would reduce the loss by moving the
primitive --- a signature of under-reconstruction. Where $\tau_i$ exceeds a
threshold, the model is refined in one of two ways:
\begin{itemize}
  \item \textbf{Cloning} (primitive smaller than a scene-scale-relative bound):
        duplicate it and displace along the positional gradient, adding capacity
        to an under-populated region;
  \item \textbf{Splitting} (primitive larger than that bound): replace it by two
        smaller primitives sampled from its own distribution, with scales divided
        by a constant, refining an over-large blob.
\end{itemize}
Periodically all opacities are reset toward zero; primitives that are genuinely
supported by the data recover, while those that merely accumulated as
semi-transparent clutter fall below the pruning threshold and are removed.

Two scene-scale dependencies deserve emphasis, because both caused real failures
in this project. First, the positional learning rate is scaled by the scene
extent; second, the clone-versus-split decision compares primitive size against a
scene-scale-relative bound. Estimating the scene extent from the raw point cloud
makes it hostage to SfM outliers: a single stray point inflates the extent, which
simultaneously raises the positional learning rate (primitives jitter, producing
blur) and raises the split threshold (primitives are only ever cloned, never
shrink). We therefore define the extent from \emph{camera centres} rather than
points, following the original spatial-scale convention.

%---------------------------------------------------------------------------
\subsection{The sparse-view problem and monocular depth priors}
\label{sec:sparse}

Equation~\eqref{eq:compositing} is severely underdetermined when views are few or
share little parallax. A set of primitives can reproduce every training image
exactly while placing mass at physically wrong depths, because depth is only
constrained by \emph{disagreement} between views. With a single-sided capture the
disagreement is small, and the photometric loss alone cannot distinguish a correct
surface from a family of depth-shifted alternatives.

The established remedy is to import geometric knowledge from a monocular depth
network \cite{ranftl2022midas,yang2024depthanythingv2} and use it to regularise
the rendered depth, as in FSGS \cite{zhu2024fsgs}, SparseGS
\cite{xiong2024sparsegs} and DNGaussian \cite{li2024dngaussian}. The rendered
depth is obtained by compositing depth exactly as colour is composited in
\eqref{eq:compositing}:
\begin{equation}
  \hat{D}(\bm{p}) \;=\; \sum_i d_i\,\hat{\alpha}_i(\bm{p}) \prod_{j<i}\big(1-\hat{\alpha}_j(\bm{p})\big).
\end{equation}

\paragraph{Why the loss must be affine-invariant.}
Monocular networks trained on mixed datasets predict depth only \emph{up to an
unknown scale and shift} \cite{ranftl2022midas}; the prediction lives in a
disparity-like space with no metric meaning. Penalising
$\lVert \hat{D}-D_{\text{mono}}\rVert$ directly would therefore impose an
arbitrary global scale on the reconstruction and fight the SfM geometry. The loss
must be invariant to the ambiguity, a requirement dating to the scale-invariant
objective of Eigen et al.~\cite{eigen2014depth}.

We adopt the Pearson correlation \cite{pearson1895} between rendered and predicted
depth over the pixels of a view,
\begin{equation}
  \mathcal{L}_{\text{depth}} \;=\; 1-\rho\big(\hat{D},\,D_{\text{mono}}\big),
  \qquad
  \rho(X,Y)=\frac{\operatorname{Cov}(X,Y)}{\sigma_X\,\sigma_Y}.
  \label{eq:pearson}
\end{equation}
For any $a>0$ and $b\in\mathbb{R}$ we have $\rho(aX+b,Y)=\rho(X,Y)$, so
\eqref{eq:pearson} is exactly invariant to the affine ambiguity: it constrains the
depth \emph{ordering} of the scene while leaving scale and offset to be fixed by
the cameras. Note the sign convention matters --- since the network emits a
disparity-like quantity (large means \emph{near}), it must be inverted before use,
or the loss would drive the geometry to the reflection of the correct surface.

The prior is ramped in only after the photometric term has established coarse
structure, so that a noisy monocular estimate does not dominate early
optimisation.

%---------------------------------------------------------------------------
\subsection{Floaters: spatial bounding and opacity/scale control}
\label{sec:floaters}

``Floaters'' are primitives occupying empty space, typically faint and often
large, which reproduce training views by acting as a diffuse haze near the cameras
while corrupting novel views. As noted after \eqref{eq:compositing}, the
compositing model gives no intrinsic preference for compact opaque geometry over
diffuse semi-transparent clutter; in the sparse-view regime such clutter is an
easy local minimum. In NeRF-style methods the standard countermeasure is the
distortion loss of Mip-NeRF~360 \cite{barron2022mipnerf360}, which penalises
weight dispersed along a ray and thereby encourages compact ray termination;
Mip-Splatting \cite{yu2024mipsplatting} addresses the related scale-aliasing
artefacts that appear when primitives are observed at sampling rates far from
those seen in training.

We use two direct, geometrically motivated constraints. First, a \emph{scene
bound}: the capture took place in a room, so admissible geometry lies within the
axis-aligned bounding box of the camera centres and the SfM points, dilated by a
margin $m$,
\begin{equation}
  \mathcal{B} \;=\; \big[\bm{l}-m(\bm{h}-\bm{l}),\;\bm{h}+m(\bm{h}-\bm{l})\big],
  \qquad
  \bm{l}=\min(\mathcal{P}\cup\mathcal{C}),\;\;
  \bm{h}=\max(\mathcal{P}\cup\mathcal{C}).
\end{equation}
Primitives leaving $\mathcal{B}$ have their opacity driven to zero and are then
removed by the existing opacity prune. Second, an explicit \emph{scale and opacity
criterion}: primitives whose largest scale exceeds a fraction of the scene extent,
or whose opacity falls below a threshold, are suppressed and finally pruned.

Suppressing rather than deleting mid-training is deliberate: the density-control
state of \S\ref{sec:densification} is indexed per primitive, so removing entries
externally would desynchronise it. Driving opacity to $\sigma(-20)\approx 2\times10^{-9}$
changes no tensor shape, makes the primitive invisible, and lets the existing
prune reclaim it.

Empirically these constraints removed $\approx\!18\text{--}20\%$ of primitives.
Crucially, the diagnostic revealed which mechanism did the work: only $0.01\%$ of
baseline primitives lay outside $\mathcal{B}$, so the bound acted as a safety net,
whereas $18.2\%$ had opacity below $0.02$ and the largest scales spanned the whole
scene. The floaters here were in-volume haze rather than runaway geometry, and it
was the opacity/scale criterion that eliminated them --- median opacity rose from
$0.51$ to $0.88$.

%---------------------------------------------------------------------------
\subsection{Surface-aligned variants}
\label{sec:2dgs}

Because 3D Gaussians are volumetric, their zero-crossing does not define a
surface, and depth composited through a thick primitive is biased. 2D Gaussian
Splatting \cite{huang20242dgs} replaces each ellipsoid with an oriented planar
disk, so the primitive \emph{is} a surface element with a well-defined normal;
combined with normal-consistency and distortion regularisers it yields markedly
better geometry, and SuGaR \cite{guedon2024sugar} similarly binds Gaussians to an
extractable mesh.

Either representation supports mesh extraction by fusing rendered depth maps into
a truncated signed-distance volume \cite{curless1996tsdf} and running marching
cubes, which is attractive here because the object of interest is a relief surface
rather than a radiance field. Note that this route does not strictly require
surface-aligned primitives: composited depth from a 3D Gaussian model can be fused
in the same way, at the cost of the depth bias induced by volumetric primitives.

We evaluated 2DGS and report a negative result: on this capture it produced flat,
featureless renderings at PSNR~$\approx13$. The reason is instructive rather than
incidental. Surface-aligned primitives must recover an orientation as well as a
position, and disk normals are constrained mainly by observing a surface from
multiple directions. A single-sided capture supplies almost no such signal, so the
added constraint removes degrees of freedom the photometric loss was using without
supplying the multi-view evidence needed to determine the ones that remain. The
same regularisation that improves well-covered scenes is therefore harmful here.

%---------------------------------------------------------------------------
\subsection{Per-image appearance modelling}
\label{sec:appearance}

Coverage is the binding constraint (\S\ref{sec:sfm}), and the cheapest source of
additional coverage is additional clips of the same object. This is obstructed by
photometry rather than geometry: a handheld camera re-runs auto-exposure and
auto-white-balance per clip, so the same surface is recorded with different
responses. The model in \eqref{eq:compositing} has no parameter expressing ``same
radiance, different camera response'', so the photometric loss
\eqref{eq:photo} cannot be driven to zero anywhere; residual gradients stay high
across the whole image, and \S\ref{sec:densification} interprets uniformly high
$\tau_i$ as under-reconstruction and densifies without bound. Our three-clip merge
duly plateaued near PSNR~$14$ while a single clip reached~$24$.

NeRF-in-the-Wild \cite{martinbrualla2021nerfw} solves this by attaching a latent
appearance code $\bm{\ell}_k$ to each image $k$ and conditioning the emitted
radiance on it. The codes are free parameters optimised jointly with the scene ---
generative latent optimisation \cite{bojanowski2018glo} --- rather than the output
of an encoder. We use the same idea with a deliberately low-capacity decoder: a
small MLP maps $\bm{\ell}_k$ to a global affine colour transform applied to the
rendered image,
\begin{equation}
  \hat{\bm{I}}_k^{\text{corr}}(\bm{p}) \;=\; \bm{M}(\bm{\ell}_k)\,\hat{\bm{I}}_k(\bm{p}) \;+\; \bm{b}(\bm{\ell}_k),
  \qquad
  \bm{M}\in\mathbb{R}^{3\times3},\;\bm{b}\in\mathbb{R}^{3},
  \label{eq:appearance}
\end{equation}
with the decoder initialised so that $\bm{M}=\bm{I}$ and $\bm{b}=\bm{0}$;
training therefore starts exactly at the identity and departs from it only if
doing so reduces the loss.

\paragraph{Why this cannot absorb geometry.}
Capacity control is the crux: a per-image correction expressive enough to repaint
pixels individually would simply memorise each training view, destroying the
multi-view constraint that makes reconstruction possible. The transform
\eqref{eq:appearance} is applied with the \emph{same} $(\bm{M},\bm{b})$ at every
pixel, so for any spatial permutation $\pi$ of the pixels,
\begin{equation}
  \bm{M}\,\hat{\bm{I}}\big(\pi(\bm{p})\big)+\bm{b}
  \;=\;
  \Big(\bm{M}\,\hat{\bm{I}}(\bm{p})+\bm{b}\Big)\circ\pi ,
\end{equation}
i.e.\ the correction commutes with any rearrangement of image content. A map that
is invariant to where things are cannot encode where things are: it carries
exactly zero spatial information and can represent only a global photometric
change. It can absorb exposure, gain and white-balance drift, and nothing else.

At evaluation time we apply the transform decoded from the \emph{mean} training
latent, giving one canonical appearance. Fitting a held-out image's own latent
would leak information from the test image into the prediction and inflate the
reported metrics.

%---------------------------------------------------------------------------
\subsection{Pose refinement}
\label{sec:pose}

A residual error source is the poses themselves: \S\ref{sec:sfm} treats them as
fixed, yet low-overlap SfM yields noisier extrinsics, and a pose error is
indistinguishable from a geometry error to the photometric loss. Joint
optimisation of poses with the scene, as in BARF \cite{lin2021barf}, addresses
this, with the caveat that a high-frequency representation admits many poor local
minima --- BARF's central observation is that the optimisation must be coarse-to-fine.
The analogous mechanism in our setting is the progressive introduction of
spherical-harmonic bands (\S\ref{sec:3dgs}).

%---------------------------------------------------------------------------
\subsection{Evaluation}
\label{sec:eval}

We report PSNR, SSIM \cite{wang2004ssim} and LPIPS \cite{zhang2018lpips} on
held-out views. The three are deliberately complementary: PSNR derives from mean
squared error and is dominated by low-frequency agreement, rewarding a blurred
prediction; SSIM compares local luminance, contrast and structure statistics and
so is sensitive to the relief we care about; LPIPS compares deep features and
correlates better with human judgement of perceptual similarity.

Two methodological cautions apply to the comparisons in this report. First, model
selection must use validation views only --- selecting on the test split silently
converts it into a validation split. Second, these metrics are \emph{not}
comparable across reconstructions built at different image resolutions or with
different held-out splits: PSNR in particular shifts with the spatial frequency
content of the images. Where we compare a higher-resolution reconstruction against
a lower-resolution one we therefore compare \emph{distributions} over views ---
medians and the frequency of catastrophic views --- and state the caveat
explicitly, rather than reporting a single number as if it were a like-for-like
improvement.

Finally, a scalar mean conceals the failure structure that matters. Per-view
analysis of our reconstruction showed the errors were not uniformly distributed
but concentrated entirely in extreme close-ups at grazing incidence, which
identified a spatial-frequency deficit --- the source images had been downscaled
--- and directly motivated the higher-resolution reconstruction. Aggregate metrics
diagnose nothing; the distribution, and the images themselves, do.

% !TeX root = main.tex
%
% Experimental setup and results for the Pavillon reconstruction.

\section{Experiments}
\label{sec:experiments}

%---------------------------------------------------------------------------
\subsection{Capture and pipeline configuration}
\label{sec:setup}

The subject is a carved wooden ceiling panel, stood vertical and filmed handheld in
a room with an iPhone (3840$\times$2160, 30\,fps, 196\,s). The capture is
single-sided --- the camera sweeps across the panel but never orbits it --- so
parallax between views is small, and the frames alternate between contextual
mid-range shots and extreme close-ups at grazing incidence.

Table~\ref{tab:settings} lists the configuration of the recommended model. Every
value is set in one YAML file; no code changes are required to reproduce it.

\begin{table}[h]
\centering\small
\begin{tabular}{llp{7.2cm}}
\toprule
Stage & Setting & Value / rationale \\
\midrule
Frames & sampling & 4\,fps, capped at 400 frames \\
       & resolution & long edge \textbf{2560}\,px (from 3840; see \S\ref{sec:res}) \\
       & blur filter & variance-of-Laplacian $\geq 50$, measured at a \emph{fixed
                       1600\,px reference scale} so the threshold is
                       resolution-independent \\
       &             & $\rightarrow$ 282 of 400 frames accepted \\
\midrule
SfM    & matcher & exhaustive (low overlap $\Rightarrow$ all pairs) \\
       & mapper  & \textbf{GLOMAP} (global) rather than incremental COLMAP \\
       & result  & \textbf{282/282} registered, 0.985\,px mean reprojection error,
                   152\,792 sparse points \\
\midrule
Split  & strategy & pose-aware; 226 train / 28 val / 28 test \\
\midrule
Train  & backend & \texttt{gsplat}, 30\,k iterations, SH degree 3 \\
       & \textbf{capacity} & \texttt{cap\_max} $=$ \textbf{375\,k} (see \S\ref{sec:capacity}) \\
       & densification & grad threshold $3\times10^{-4}$, interval 100,
                         iterations 500--15\,000, opacity reset every 3000 \\
       & room bounds & AABB(cameras $\cup$ points) $\times1.5$, enforced every 500 steps \\
       & anti-floater & suppress scale $>0.15\cdot$extent; final prune below
                        opacity 0.005 \\
       & depth prior & DepthAnything-v2-Small, Pearson loss, weight 0.1, from iter 2000 \\
\bottomrule
\end{tabular}
\caption{Configuration of the recommended model (\texttt{pavillon\_orbit\_hidetail\_cap375k}).}
\label{tab:settings}
\end{table}

%---------------------------------------------------------------------------
\subsection{Results}

\begin{table}[h]
\centering\small
\begin{tabular}{lrrrrrr}
\toprule
Model & Gaussians & PSNR & median & SSIM & LPIPS & $<$18\,dB \\
\midrule
\multicolumn{7}{l}{\emph{1600\,px dataset, 18 test views}}\\
Baseline (GLOMAP + 3DGS)      & 1\,533\,137 & 23.89 & --- & 0.8249 & 0.2451 & --- \\
A3 regularized                 & 1\,216\,252 & 22.29 & 23.43 & 0.7958 & 0.2901 & 17\,\% \\
A3 ablation (depth off)        & 1\,232\,607 & 22.55 & 22.76 & 0.7978 & 0.2941 & 17\,\% \\
2DGS backend                   &   277\,540 & 13.31 & --- & 0.5778 & 0.8473 & --- \\
\midrule
\multicolumn{7}{l}{\emph{2560\,px dataset, 28 test views}}\\
cap 1.5\,M, ADC                & 1\,244\,835 & 23.08 & 24.48 & 0.8461 & 0.3096 & 11\,\% \\
cap 750\,k, ADC                &   653\,831 & 24.32 & 25.10 & 0.8596 & \textbf{0.3025} & 11\,\% \\
\textbf{cap 375\,k, ADC (recommended)} & \textbf{358\,878} & \textbf{24.92} & 25.13 & \textbf{0.8619} & 0.3132 & \textbf{0\,\%} \\
cap 190\,k, ADC                &   192\,475 & 24.86 & \textbf{25.18} & 0.8559 & 0.3382 & \textbf{0\,\%} \\
$+$ pose refinement            &   354\,250 & 23.94 & 24.40 & 0.8290 & 0.3301 & 0\,\% \\
cap 375\,k, \textbf{MCMC}      &   365\,497 & 24.83 & \textbf{25.32} & \textbf{0.8623} & \textbf{0.3094} & 0\,\% \\
\midrule
\multicolumn{7}{l}{\emph{multi-clip dataset (2 clips), 33 test views}}\\
cap 3.0\,M                     & 2\,502\,172 & 21.36 & 21.23 & 0.8243 & 0.3460 & 33\,\% \\
cap 1.5\,M                     & 1\,225\,089 & 22.03 & 22.82 & 0.8332 & 0.3420 & 24\,\% \\
cap 750\,k                     &   637\,412 & 22.37 & 23.52 & 0.8396 & 0.3350 & 21\,\% \\
\bottomrule
\end{tabular}
\caption{All trained models. \textbf{Rows in different blocks are not directly
comparable}: they are separate datasets with their own SfM solution, held-out split
and image resolution, and PSNR/LPIPS shift with resolution. Comparisons within a
block are valid.}
\label{tab:results}
\end{table}

\paragraph{Resolution.}
\label{sec:res}
Per-view analysis of the 1600\,px model showed the errors were not spread evenly:
every catastrophic view (12.7--18.8\,dB) was an extreme close-up at grazing
incidence, while mid-range views reached 25.6--26.5\,dB. That is a
spatial-frequency deficit --- the pipeline was downscaling a 4K source to 1600\,px,
discarding exactly the carving detail those views demand. Re-running at 2560\,px
with denser frame sampling raised registration to 282/282 (from 181/193) and lifted
median PSNR from 23.43 to 24.48 with SSIM 0.796\,$\rightarrow$\,0.846. The aggregate
mean had concealed this entirely; only the per-view distribution exposed it.

\paragraph{Capacity.}
\label{sec:capacity}
The largest late gain came from making the model \emph{smaller}. Holding everything
else fixed and varying only \texttt{cap\_max}, PSNR and SSIM peak at 375\,k while
the catastrophic-view tail disappears below 750\,k. Cutting the budget $4\times$
from 1.5\,M gained 1.8\,dB and shrank the model $3.5\times$ (295\,MB
$\rightarrow$ 89\,MB). This is sparse-view overfitting: with low parallax the
objective is underdetermined, so surplus primitives buy configurations that fit the
training views and fail on held-out ones (\S\ref{sec:sparse}). The curve genuinely
turns --- 190\,k is worse than 375\,k on PSNR and SSIM --- so 375\,k is an optimum
rather than the smallest budget tried.

The metrics disagree informatively: \textbf{LPIPS is minimised at 750\,k} and
degrades monotonically as capacity shrinks, while PSNR, SSIM and robustness prefer
less. Perceptual detail wants capacity; accuracy and generalisation want less. We
report both operating points rather than collapsing them into one recommendation.

\paragraph{Ablations and negative results.}
\begin{itemize}
  \item \textbf{Depth prior:} disabling it changes test PSNR by only 0.26\,dB
        (22.29\,$\rightarrow$\,22.55). It is nearly free and slightly \emph{increases}
        floater removal, but it is not what determines quality here. The
        $\approx$1.4\,dB cost of the A3 configuration is the final opacity prune, not
        the prior --- an ablation that corrected our initial assumption.
  \item \textbf{Anti-floater:} removes 18--20\,\% of primitives. Baseline models
        carried \textbf{18.2\,\%} of Gaussians below 0.02 opacity with scales up to
        1.0 (scene-spanning); the regularized model has \textbf{0\,\%} sub-0.02 haze,
        max scale 0.16, and median opacity 0.51\,$\rightarrow$\,0.88.
  \item \textbf{2DGS:} collapses to flat renders (13.3\,dB). Disk normals are
        constrained by observing a surface from multiple directions, which a
        single-sided capture does not provide.
  \item \textbf{Multi-clip merging:} SfM merges the clips cleanly (332/332 at
        0.921\,px, 266 train views) and the appearance latents demonstrably absorb
        exposure drift, yet the merged model still trails the best single-clip model
        by $\approx$1\,dB. The second clip enlarges the reconstructed volume while
        adding few views.
  \item \textbf{Pose refinement:} costs 1\,dB. The learned deltas moved only
        $0.036^\circ$ and $7\times10^{-4}$ --- about one pixel --- because GLOMAP's
        poses were already sub-pixel accurate. Refining only training cameras (test
        poses must not be optimised) then misaligns the scene slightly relative to
        the evaluation cameras.
\end{itemize}

\paragraph{Densification strategy.}
Replacing the heuristic adaptive density control with 3DGS-MCMC
\cite{kheradmand2024mcmc} at the same 375\,k budget changes nothing measurable:
the paired per-view difference is $-0.09 \pm 0.45$\,dB (95\,\% CI over 28 held-out
views), i.e.\ indistinguishable. MCMC is ahead on 17 of 28 views and on median,
SSIM and LPIPS; the heuristic wins the mean only through a single $-5.32$\,dB
outlier. We therefore make no quality claim in either direction --- with 28 test
views the experiment cannot resolve differences below roughly half a decibel. The
result is nonetheless informative: it indicates the large gain in
\S\ref{sec:capacity} came from the \emph{budget}, not from how the budget is
allocated, which is what the sparse-view overfitting account predicts. MCMC remains
preferable operationally, since it eliminates the gradient threshold and makes the
budget an explicit target rather than an emergent ceiling.

\paragraph{Two measurement bugs worth reporting.}
Both would have silently corrupted conclusions. First, the blur filter scored
variance-of-Laplacian on the full-resolution frame, a resolution-dependent
statistic: raising the extraction resolution to 2560\,px pushed the median score
from 92 to 36 and rejected 308/400 frames as ``blurry'', which would have made the
higher-resolution run \emph{worse} than the one it replaced. Second, evaluation
rendered without restoring the appearance model, scoring held-out views on the
uncorrected output while training optimised corrected ones; fixing it moved the
multi-clip model from 19.0 to 22.0\,dB. Neither was visible in an aggregate number
--- the first surfaced from a frame-count sanity check, the second from noticing that
PSNR had collapsed while SSIM barely moved, which is the signature of a global
photometric offset rather than broken geometry.

%---------------------------------------------------------------------------
\subsection{Geometry output}

Fusing the composited depth of the recommended model over all 226 training cameras
into a TSDF volume \cite{curless1996tsdf} yields a mesh of 952\,846 vertices and
1\,842\,626 triangles, cropped to the object bounding box. The panel's planks and
carved cross-members are resolved. Two limitations are inherent: volumetric
primitives give a biased expected depth, so the surface is softer than a
surface-aligned method would produce, and the fusion covers whatever the cameras
saw, so the crop is what isolates the subject from the room.


\bibliographystyle{plainnat}
\bibliography{references}

\end{document}
